\documentclass[12pt]{extreport}
\usepackage[utf8]{inputenc}

% --- 1. FONTS & FORMATTING ---
\usepackage{mathptmx}  % Times New Roman
\usepackage[T1]{fontenc}
\usepackage{scrextend}
\changefontsizes{13pt}

% --- 2. MARGINS ---
\usepackage[left=2.5cm, top=2.0cm, right=1.5cm, bottom=2.5cm]{geometry}

% --- 3. SPACING ---
\usepackage{setspace}
\onehalfspacing

% --- 4. PACKAGES ---
\usepackage{graphicx}
\usepackage{titlesec}
\usepackage{indentfirst}
\usepackage{float}
\usepackage{hyperref}
\usepackage{listings}
\usepackage{xcolor}
\usepackage{amsmath}

% --- 5. HEADING FORMAT ---
\titleformat{\chapter}[display]
  {\normalfont\bfseries\centering}
  {\MakeUppercase{\chaptertitlename}\ \thechapter}
  {0pt}
  {\MakeUppercase}
\titlespacing*{\chapter}{0pt}{-20pt}{24pt}

% --- 6. CODE LISTING STYLE ---
\lstdefinestyle{tsstyle}{
    language=Java, % closest built-in; we add TS keywords below
    basicstyle=\ttfamily\small,
    keywordstyle=\color{blue}\bfseries,
    commentstyle=\color{gray}\itshape,
    stringstyle=\color{red},
    numbers=left,
    numberstyle=\tiny\color{gray},
    numbersep=8pt,
    frame=single,
    breaklines=true,
    tabsize=2,
    showstringspaces=false,
    captionpos=b,
    morekeywords={const,let,function,export,interface,import,from,type,number,boolean,string,null,undefined,as,of}
}

\begin{document}

% --- TITLE PAGE ---
\begin{titlepage}
    \centering
    \vspace*{1cm}
    \textbf{VIETNAM NATIONAL UNIVERSITY -- HO CHI MINH CITY}\\
    \textbf{INTERNATIONAL UNIVERSITY}\\
    \textbf{SCHOOL OF COMPUTER SCIENCE AND ENGINEERING}

    \vspace{3cm}
    \textbf{ADVANCED COMPUTER GRAPHIC}\\
    \textbf{ASSIGNMENT 2}

    \vspace{4cm}
    \textbf{ASSIGNMENT REPORT}

    \vspace{3cm}
    \begin{flushleft}
    \hspace{4cm} Student: Tran Bao Thanh, MITIU25214 \\
    \end{flushleft}

    \vfill
    March 2026
\end{titlepage}

% --- TABLE OF CONTENTS ---
\pagenumbering{roman}
\setcounter{page}{2}
\tableofcontents
\newpage

% ===========================================================================
\pagenumbering{arabic}
\chapter{FILLING HOLES ON THE 2D SURFACE OF A POINT CLOUD}
% ===========================================================================

When we have a 2D point cloud with gaps or sparse regions, we need a way to fill those holes with a good quality triangle mesh. The approach we took is to first build a Delaunay triangulation of the point cloud using the Bowyer-Watson algorithm, then identify triangles that are too large or too thin, these are the ones covering the holes and sparse areas, and then use Ruppert's refinement algorithm to insert new points into those bad triangles until the mesh is filled in properly.

We built a web application to demonstrate this process using Vue 3 (Composition API) and Vite. The core algorithms which are Bowyer-Watson, circumcircle computation, and Ruppert's refinement we wrote from scratch in TypeScript. The convex hull uses the \texttt{d3-polygon} library since that is not part of the main algorithms we are studying in this assignment.

The application lets you generate a random 2D point cloud, then it triangulates the points and highlights where the holes and bad triangles are. You can then refine the mesh step by step or all at once and watch as the holes get filled in with new points and better shaped triangles.

\section{Project Structure}

\begin{verbatim}
src/
    geometry.ts       -- Types + math primitives (circumcircle, area, angle)
    triangulation.ts  -- Bowyer-Watson algorithm
    refinement.ts     -- Ruppert's refinement
    hull.ts           -- Convex hull (d3-polygon wrapper)
    App.vue           -- State, UI, rendering
\end{verbatim}

\section{Data Types}

We defined these types in \texttt{geometry.ts} and they are shared across all the other files:

\begin{lstlisting}[style=tsstyle, caption={Shared types (\texttt{geometry.ts})}]
interface Point {
  x: number
  y: number
  __super?: boolean  // marks super-triangle vertices
}

interface Circle {
  x: number
  y: number
  r: number
}

interface TriangleInfo {
  p1: Point; p2: Point; p3: Point
  area: number
  angle: number        // smallest interior angle (degrees)
  cc: Circle | null    // circumscribed circle (null if degenerate)
  bigArea: boolean     // area exceeds threshold
  smallAngle: boolean  // min angle below threshold
  bad: boolean         // needs refinement (bigArea OR smallAngle)
}
\end{lstlisting}

We also have some small utility functions that get used everywhere:

\begin{lstlisting}[style=tsstyle, caption={Utility functions (\texttt{geometry.ts})}]
const EPSILON = 1e-10  // tolerance for float comparisons

// prevent NaN/Infinity from crashing SVG rendering
function safe(v: number): number {
  return isFinite(v) ? v : 0
}

// needed before passing to Math.acos
function clamp(value: number): number {
  return Math.max(-1, Math.min(1, value))
}

function distanceSq(a: Point, b: Point): number {
  return (a.x - b.x) ** 2 + (a.y - b.y) ** 2
}

function distance(a: Point, b: Point): number {
  return Math.sqrt(distanceSq(a, b))
}

function midpoint(a: Point, b: Point): Point {
  return { x: (a.x + b.x) / 2, y: (a.y + b.y) / 2 }
}
\end{lstlisting}

\section{Core Algorithm: Bowyer-Watson}

Bowyer-Watson is an incremental algorithm where each point in the point set is inserted one at a time into an existing triangulation. What the algorithm does is after every insertion it checks whether any of the existing triangles have their circumcircle violated by the new point. If a point lands inside the circumcircle of a triangle then that triangle is no longer valid under the Delaunay condition so it has to be removed. The algorithm then takes the hole that is left behind from removing those triangles and fills it by connecting the edges of the hole back to the newly inserted point.

\begin{enumerate}
    \item First we need a starting triangulation so we create a very large ``super-triangle'' that is big enough to contain every single point in our input set inside of it.
    \item Then for each point we want to insert:
    \begin{enumerate}
        \item Go through all the triangles and check which ones have a circumcircle that contains this new point, those are the ones that break the Delaunay rule.
        \item Delete all of those bad triangles. This leaves a polygon shaped hole in the mesh.
        \item Look at the edges of the deleted triangles and figure out which edges are on the boundary of the hole. An edge is on the boundary if it only belonged to one deleted triangle. If two deleted triangles shared an edge then that edge is on the inside of the hole not the boundary.
        \item Take each boundary edge and make a new triangle by connecting it to the new point.
    \end{enumerate}
    \item After all the points are in, we go back and delete any triangle that still has one of the super-triangle's vertices. Those were just temporary scaffolding to get the algorithm started.
\end{enumerate}

The full implementation is in \texttt{triangulation.ts}. First the helper functions for edge detection:

\begin{lstlisting}[style=tsstyle, caption={Edge helpers (\texttt{triangulation.ts})}]
interface Triangle {
  p: [Point, Point, Point]
}

// check if two edges are the same (works in either direction)
function edgesMatch(e1: [Point, Point], e2: [Point, Point]): boolean {
  return (e1[0] === e2[1] && e1[1] === e2[0])
      || (e1[0] === e2[0] && e1[1] === e2[1])
}

// get all 3 edges of a triangle
function triangleEdges(tri: Triangle): [Point, Point][] {
  return [
    [tri.p[0], tri.p[1]],
    [tri.p[1], tri.p[2]],
    [tri.p[2], tri.p[0]]
  ]
}

// is this edge shared with another triangle in the list?
function isSharedEdge(
  edge: [Point, Point], ownerTri: Triangle, allTris: Triangle[]
): boolean {
  return allTris.some(other => {
    if (other === ownerTri) return false
    return triangleEdges(other).some(
      otherEdge => edgesMatch(edge, otherEdge)
    )
  })
}

// find boundary edges of the hole left by removing triangles
function findBoundaryEdges(
  removedTriangles: Triangle[]
): [Point, Point][] {
  const boundary: [Point, Point][] = []
  for (const tri of removedTriangles) {
    for (const edge of triangleEdges(tri)) {
      if (!isSharedEdge(edge, tri, removedTriangles)) {
        boundary.push(edge)
      }
    }
  }
  return boundary
}
\end{lstlisting}

The circumcircle containment test that Bowyer-Watson uses to decide which triangles to remove:

\begin{lstlisting}[style=tsstyle, caption={Circumcircle containment test (\texttt{triangulation.ts})}]
function isInsideCircumcircle(point: Point, tri: Triangle): boolean {
  const cc = circumcircle(tri.p[0], tri.p[1], tri.p[2])
  if (!cc) return false

  const distToCenter = (point.x - cc.x) ** 2 + (point.y - cc.y) ** 2
  const radiusSq = cc.r ** 2
  return distToCenter <= radiusSq + EPSILON
}
\end{lstlisting}

And the main Bowyer-Watson function that puts it all together:

\begin{lstlisting}[style=tsstyle, caption={Bowyer-Watson algorithm (\texttt{triangulation.ts})}]
function bowyerWatson(points: Point[]): Triangle[] {
  if (points.length < 3) return []

  // Step 1: super-triangle that contains all points
  const minX = Math.min(...points.map(p => p.x)) - 100
  const maxX = Math.max(...points.map(p => p.x)) + 100
  const minY = Math.min(...points.map(p => p.y)) - 100
  const maxY = Math.max(...points.map(p => p.y)) + 100
  const padding = Math.max(maxX - minX, maxY - minY) * 2

  const superA: Point = { x: minX - padding, y: minY - 1, __super: true }
  const superB: Point = { x: minX + padding, y: minY - 1, __super: true }
  const superC: Point = {
    x: (minX + maxX) / 2, y: maxY + padding, __super: true
  }

  let triangles: Triangle[] = [{ p: [superA, superB, superC] }]

  // Step 2: insert each point
  for (const point of points) {
    const invalidated: Triangle[] = []
    const kept: Triangle[] = []

    for (const tri of triangles) {
      if (isInsideCircumcircle(point, tri)) {
        invalidated.push(tri)
      } else {
        kept.push(tri)
      }
    }

    const boundary = findBoundaryEdges(invalidated)

    triangles = kept
    for (const [edgeA, edgeB] of boundary) {
      triangles.push({ p: [edgeA, edgeB, point] })
    }
  }

  // Step 3: remove super-triangle scaffolding
  return triangles.filter(tri => !tri.p.some(p => p.__super))
}
\end{lstlisting}

\subsection{Circumcircle Computation}

If you have three points that are not all sitting on the same line then there is always going to be exactly one circle you can draw that passes through all three of them, that circle is what we call the circumcircle. The center of it is the circumcenter. You can find the circumcenter by drawing the perpendicular bisector of two of the triangle's edges and seeing where they cross each other, the crossing point is the circumcenter. We implemented this in \texttt{geometry.ts} using line equation helpers:

\begin{lstlisting}[style=tsstyle, caption={Line geometry helpers (\texttt{geometry.ts})}]
interface Line { a: number; b: number; c: number }

// derive line equation ax + by = c from two points
function lineFromPoints(p1: Point, p2: Point): Line {
  const a = p2.y - p1.y
  const b = p1.x - p2.x
  const c = a * p1.x + b * p1.y
  return { a, b, c }
}

// rotate line 90 deg and force through midpoint
function perpendicularBisector(line: Line, mid: Point): Line {
  const a = -line.b        // swap and negate one
  const b = line.a
  const c = a * mid.x + b * mid.y  // plug in midpoint
  return { a, b, c }
}

// Cramer's rule: solve two lines for intersection
function lineIntersection(l1: Line, l2: Line): Point | null {
  const det = l1.a * l2.b - l2.a * l1.b
  if (Math.abs(det) < EPSILON) return null // parallel = collinear

  const x = (l2.b * l1.c - l1.b * l2.c) / det
  const y = (l1.a * l2.c - l2.a * l1.c) / det
  if (!isFinite(x) || !isFinite(y)) return null
  return { x, y }
}
\end{lstlisting}

\begin{lstlisting}[style=tsstyle, caption={Circumcircle computation (\texttt{geometry.ts})}]
function circumcircle(p1: Point, p2: Point, p3: Point): Circle | null {
  // Step 1: line equations for two edges
  const lineAB = lineFromPoints(p1, p2)
  const lineBC = lineFromPoints(p2, p3)

  // Step 2: perpendicular bisectors through each midpoint
  const midAB = midpoint(p1, p2)
  const midBC = midpoint(p2, p3)
  const bisectorAB = perpendicularBisector(lineAB, midAB)
  const bisectorBC = perpendicularBisector(lineBC, midBC)

  // Step 3: circumcenter = where the bisectors cross
  const center = lineIntersection(bisectorAB, bisectorBC)
  if (!center) return null // collinear points, no circle

  // Step 4: radius = distance from center to any vertex
  const radius = distance(center, p1)
  if (!isFinite(radius)) return null

  return { x: center.x, y: center.y, r: radius }
}
\end{lstlisting}

\subsection{Point-in-Triangle Test}

We need this test in a couple of places, in the refinement we use it to check if the circumcenter falls inside the triangle and we also use it for the click interaction in the UI. The way it works is it computes the cross product at each edge and then checks if they all have the same sign:

\begin{lstlisting}[style=tsstyle, caption={Point-in-triangle test (\texttt{geometry.ts})}]
function pointInTriangle(
  px: number, py: number, p1: Point, p2: Point, p3: Point
): boolean {
  const cross1 = (px - p2.x) * (p1.y - p2.y)
               - (p1.x - p2.x) * (py - p2.y)
  const cross2 = (px - p3.x) * (p2.y - p3.y)
               - (p2.x - p3.x) * (py - p3.y)
  const cross3 = (px - p1.x) * (p3.y - p1.y)
               - (p3.x - p1.x) * (py - p1.y)

  const hasNegative = (cross1 < -EPSILON)
                   || (cross2 < -EPSILON)
                   || (cross3 < -EPSILON)
  const hasPositive = (cross1 > EPSILON)
                   || (cross2 > EPSILON)
                   || (cross3 > EPSILON)

  // inside only if all cross products have the same sign
  return !(hasNegative && hasPositive)
}
\end{lstlisting}

\section{Triangle Quality Metrics}

Not every triangle that comes out of the Delaunay triangulation is going to look good. Some of them end up really large, some end up really thin and pointy. We need a way to flag these so we know which ones need to be fixed. A triangle gets flagged as ``bad'' if its area is too large, the default threshold we set is 8000 but the user can adjust it with a slider, or if its smallest angle is too small which defaults to 20\textdegree and is also adjustable. The code for both checks:

\begin{lstlisting}[style=tsstyle, caption={Triangle area via cross product (\texttt{geometry.ts})}]
function triangleArea(p1: Point, p2: Point, p3: Point): number {
  const abX = p2.x - p1.x
  const abY = p2.y - p1.y
  const acX = p3.x - p1.x
  const acY = p3.y - p1.y

  const crossProduct = abX * acY - acX * abY
  const area = Math.abs(crossProduct) / 2
  return isFinite(area) ? area : 0
}
\end{lstlisting}

\begin{lstlisting}[style=tsstyle, caption={Minimum angle via law of cosines (\texttt{geometry.ts})}]
function minAngle(p1: Point, p2: Point, p3: Point): number {
  const sideA = distance(p2, p3)  // opposite p1
  const sideB = distance(p1, p3)  // opposite p2
  const sideC = distance(p1, p2)  // opposite p3

  if (sideA < EPSILON || sideB < EPSILON || sideC < EPSILON) return 0

  const aSq = sideA ** 2, bSq = sideB ** 2, cSq = sideC ** 2

  // law of cosines: angle = acos((b^2 + c^2 - a^2) / 2bc)
  const angleAtP1 = Math.acos(clamp(
    (bSq + cSq - aSq) / (2 * sideB * sideC)))
  const angleAtP2 = Math.acos(clamp(
    (aSq + cSq - bSq) / (2 * sideA * sideC)))
  const angleAtP3 = Math.acos(clamp(
    (aSq + bSq - cSq) / (2 * sideA * sideB)))

  const minRadians = Math.min(angleAtP1, angleAtP2, angleAtP3)
  return minRadians * (180 / Math.PI)  // convert to degrees
}
\end{lstlisting}

\section{Ruppert's Refinement Algorithm}

So now we know which triangles are bad and we need to fix them, this is where Ruppert's refinement comes in. The idea is you pick a bad triangle, find a good spot inside it to put a new point, add that point to the point set, and then re-run Bowyer-Watson so the whole mesh gets re-triangulated with the new point included. The full implementation is in \texttt{refinement.ts}:

\begin{lstlisting}[style=tsstyle, caption={Helper functions for insertion point selection (\texttt{refinement.ts})}]
const MAX_INSERTIONS_PER_STEP = 12  // cap per step
const MIN_DISTANCE_SQ = 100         // 10px squared

// midpoint of the longest of the 3 edges
function longestEdgeMidpoint(
  p1: Point, p2: Point, p3: Point
): Point {
  const edges: [Point, Point][] = [[p1, p2], [p2, p3], [p1, p3]]
  let longestEdge = edges[0]
  let longestLenSq = 0

  for (const [a, b] of edges) {
    const lenSq = (a.x - b.x) ** 2 + (a.y - b.y) ** 2
    if (lenSq > longestLenSq) {
      longestLenSq = lenSq
      longestEdge = [a, b]
    }
  }
  return {
    x: (longestEdge[0].x + longestEdge[1].x) / 2,
    y: (longestEdge[0].y + longestEdge[1].y) / 2
  }
}

// average of the 3 vertex positions
function centroid(p1: Point, p2: Point, p3: Point): Point {
  return {
    x: (p1.x + p2.x + p3.x) / 3,
    y: (p1.y + p2.y + p3.y) / 3
  }
}

// is this point too close to any existing point?
function isTooClose(
  candidate: Point, existingPoints: Point[]
): boolean {
  return existingPoints.some(
    p => (p.x - candidate.x) ** 2
       + (p.y - candidate.y) ** 2 < MIN_DISTANCE_SQ
  )
}
\end{lstlisting}

\begin{lstlisting}[style=tsstyle, caption={Insertion point selection: circumcenter then fallbacks (\texttt{refinement.ts})}]
function pickInsertionPoint(tri: TriangleInfo): Point | null {
  const { p1, p2, p3, cc } = tri

  if (cc) {
    // try circumcenter first (best for angle improvement)
    const circumcenter = { x: cc.x, y: cc.y }
    if (pointInTriangle(circumcenter.x, circumcenter.y, p1, p2, p3)) {
      return circumcenter
    }
    // circumcenter outside triangle -> midpoint of longest edge
    return longestEdgeMidpoint(p1, p2, p3)
  }

  // circumcircle failed (degenerate) -> centroid
  return centroid(p1, p2, p3)
}
\end{lstlisting}

\begin{lstlisting}[style=tsstyle, caption={One refinement pass (\texttt{refinement.ts})}]
function refineOnce(
  triangleData: TriangleInfo[],
  allPoints: Point[],
  hull: Point[]
): Point[] {
  // sort bad triangles by area descending (worst first)
  const badTriangles = triangleData
    .filter(t => t.bad)
    .sort((a, b) => b.area - a.area)

  if (badTriangles.length === 0) return []

  const newPoints: Point[] = []

  for (const tri of badTriangles) {
    if (newPoints.length >= MAX_INSERTIONS_PER_STEP) break

    const candidate = pickInsertionPoint(tri)
    if (!candidate || !isFinite(candidate.x)
                   || !isFinite(candidate.y)) continue

    // must be inside original point cloud boundary
    if (!pointInConvexHull(candidate, hull)) continue

    // must not be too close to existing points
    if (isTooClose(candidate, [...allPoints, ...newPoints]))
      continue

    newPoints.push(candidate)
  }
  return newPoints
}
\end{lstlisting}

After the new points have been inserted into the point set, the Bowyer-Watson algorithm is run again from scratch on the entire point set to produce a new Delaunay triangulation that incorporates the newly added points. We have to do this because inserting points changes the topology of the mesh so the triangulation needs to be recomputed from the beginning. The ``Refine All'' button in the interface repeats this refinement process multiple times until either there are no more bad triangles left in the mesh, or the algorithm cannot insert any more points because of the distance and boundary constraints.

\section{Convex Hull}

We use \texttt{d3-polygon} for the convex hull since it is not part of the algorithms we are studying. The hull is used to prevent the refinement from adding points outside the original point cloud boundary. The wrapper is in \texttt{hull.ts}:

\begin{lstlisting}[style=tsstyle, caption={Convex hull wrapper (\texttt{hull.ts})}]
import { polygonHull, polygonContains } from 'd3-polygon'

function convexHull(points: Point[]): Point[] {
  const coords: [number, number][] = points.map(p => [p.x, p.y])
  const hull = polygonHull(coords)
  if (!hull) return points.slice(0, 3)
  return hull.map(([x, y]) => ({ x, y }))
}

function pointInConvexHull(point: Point, hull: Point[]): boolean {
  if (hull.length < 3) return true
  const polygon: [number, number][] = hull.map(p => [p.x, p.y])
  return polygonContains(polygon, [point.x, point.y])
}
\end{lstlisting}

\section{User Interface}

For the front end visualization, the application renders all of the geometry inside a single SVG element with a viewBox of 700$\times$600 pixels. The user interacts with the application through controls that are placed above the SVG canvas:

\begin{itemize}
    \item \textbf{Random} --- generate 28 random points and triangulate.
    \item \textbf{Refine Step} --- run one refinement pass.
    \item \textbf{Refine All} --- refine until no bad triangles remain.
    \item \textbf{Reset} --- remove all refinement-added points.
    \item \textbf{Circumcircles} --- toggle circumcircle overlay (dashed circles).
    \item \textbf{Max Area / Min Angle sliders} --- adjust quality thresholds.
    \item \textbf{Click on canvas} --- add a point manually.
    \item \textbf{Click on triangle} --- select it and view its stats.
\end{itemize}

We use color to communicate triangle quality to the user. Each triangle in the SVG canvas is filled with a color that corresponds to whether it passed or failed the quality thresholds. We chose to do it this way because it makes it immediately obvious where the problem areas are in the mesh without having to click on individual triangles:
\begin{itemize}
    \item \textbf{Green} --- good (passes both thresholds)
    \item \textbf{Orange} --- area too large
    \item \textbf{Yellow} --- minimum angle too small
    \item \textbf{Red} --- both problems
    \item \textbf{Purple} --- currently selected
\end{itemize}

For the points we also use color to differentiate between the original random points and the ones that were inserted during refinement. The original points show up as dark circles and the refinement inserted points show up as purple circles. We thought this distinction was important because it lets you see how much work the refinement algorithm actually did compared to the starting point set. There is also a statistics bar at the bottom of the page that shows total points, how many were added by refinement, triangle count, bad triangle count, and the average minimum angle across the whole mesh.

\section{Results}

To test the application we generated a set of 28 random points scattered across the canvas. The initial Delaunay triangulation of these points contained a large number of bad triangles, shown in orange and yellow in Figure \ref{fig:web_before}. After clicking the ``Refine All'' button the refinement algorithm ran multiple passes and inserted additional points into the mesh until all the triangles passed both the area and minimum angle quality thresholds. The result after refinement is shown in Figure \ref{fig:web_after}, most of the triangles are green now which indicates they are within acceptable quality. Figure \ref{fig:web_circumcircles} shows the circumcircle overlay, the dashed blue circles represent each triangle's circumcircle and we can verify that no point falls inside another triangle's circumcircle which confirms that the Delaunay property is maintained.

\begin{figure}[H]
    \centering
    \includegraphics[width=0.85\textwidth]{web_before_refine.png}
    \caption{Initial Delaunay triangulation of 28 random points. Orange and yellow triangles are ``bad'' (too large or too skinny).}
    \label{fig:web_before}
\end{figure}

\begin{figure}[H]
    \centering
    \includegraphics[width=0.85\textwidth]{web_after_refine.png}
    \caption{After refinement. Purple dots are inserted points. Most triangles are now green (pass both quality thresholds).}
    \label{fig:web_after}
\end{figure}

\begin{figure}[H]
    \centering
    \includegraphics[width=0.85\textwidth]{web_circumcircles.png}
    \caption{Circumcircle overlay enabled. Dashed blue circles show each triangle's circumscribed circle. No point falls inside another triangle's circumcircle, confirming the Delaunay property.}
    \label{fig:web_circumcircles}
\end{figure}

\end{document}
